\section{The counting problem and its normalization}
\label{sec:intro}
We count subgroups as subsets of $S_n$, rather than conjugacy classes
or abstract isomorphism classes. All unspecified logarithms in counting
estimates are to base two; $\log_e$ denotes the natural logarithm.
Roney-Dougal and Tracey~\cite{RoneyDougalTracey2025} prove
\begin{equation}\label{eq:coarse}
 s_n\leq 2^{n^2/16+O(n^{3/2})}.
\end{equation}
Together with the elementary binary construction, this identifies the
quadratic exponential scale. The more precise benchmark considered here
also records the labelled orbit decomposition.

Put $R=\lfloor n/2\rfloor$ and $\eps=n-2R$, and define
\begin{equation}\label{eq:normalization}
 P(y)=\frac y2+\frac{y^2}6+\frac{y^4}{384},\qquad
 c_R=[y^R]e^{P(y)},\qquad
 p_n=c_R+\frac{\eps}{6}c_{R-1},\qquad L_n=n!G_Rp_n.
\end{equation}
Here $c_j=0$ for $j<0$, $G_R=\sum_k\gauss Rk2$, and $G_0=1$.
In particular $p_n$ is a coefficient, whereas the normalized number of
subgroups in the critical family will be denoted $C_n^{\crit}$.
The identity $P'(y)e^{P(y)}=(e^{P(y)})'$ gives
\begin{equation}\label{eq:coefficient-recurrence}
 R c_R=\frac12c_{R-1}+\frac13c_{R-2}+\frac1{96}c_{R-4}.
\end{equation}
All coefficients are positive. Useful crude consequences include
$c_{R-1}/c_R\leq2R$ and
$c_{R-h}/c_R\leq(2R)^h$ for $0\leq h\leq R$.

\begin{theorem}[Critical-family asymptotic]\label{thm:critical-intro}
Consider the subgroups whose even orbits are the four actions in
Table~\ref{tab:critical}. In odd degree allow, in addition, either one
fixed point or one natural $S_3$ orbit. Require full projection onto
each displayed action, and include every subgroup with these orbit
projections. Their number is $L_n(1+O(2^{-a n}))$ for some $a>0$.
Consequently $s_n\geq L_n(1-O(2^{-a n}))$.
\end{theorem}

\begin{theorem}[Section capacity]\label{thm:capacity-intro}
Let $T\leq S_s$ be faithful, transitive, and not a $2$-group, where
$s\geq24$ is even. Let $A$ be any section of its binary permutation
module, with action factoring through a specified quotient $B$ of $T$.
For a finite $\F_2B$-module $M$, put
\[
 r(M)=\max_X\frac{m_X(\Soc M)}{\dim_{\End_BX}X},\qquad r(0)=0,
\]
where $X$ runs over simple modules. There is a subspace $C\leq A^B$
such that
\begin{equation}\label{eq:capacity-intro}
 2\dim C+4r(A/C)\leq\frac{47s}{48}.
\end{equation}
The statement concerns the actual quotient $A/C$, including all socle
constituents exposed by the quotient.
\end{theorem}

The critical model and the capacity estimate are proved in
Sections~\ref{sec:critical} and~\ref{sec:capacity}. The complete counting
argument gives the following theorem.

\begin{theorem}[Precise subgroup asymptotic]\label{thm:main}
There exist $c,C>0$ and $n_0$ such that
\[
 \left|\frac{s_n}{L_n}-1\right|\leq C2^{-cn}\qquad(n\geq n_0).
\]
The constants are common to both parities. No boundedness assumption
on $s_n/L_n$ is needed.
\end{theorem}

We first prove complete counting estimates on one complementary source,
with the normalizer of each original permutation action retained.
Sections~\ref{sec:nonbinary-alphabet} and~\ref{sec:binary-complete}
give the exhaustive nonbinary action reduction and the independent binary
recurrence. The natural $C_3/A_4$ packets and exact marker collapse complete
the disjoint partition in Section~\ref{sec:assembly}. Two successive
boundedness inductions then prove Theorem~\ref{thm:main}.
The split $c=1$ calculation in Section~\ref{sec:first-c1} is a required
application of the joint-source method, with its shared-character and
inverse-complement controls retained. Section~\ref{sec:analysis} evaluates
$L_n$ and gives the saddle-point and elementary forms of the theorem.

\subsection{Actual actions and recovered profiles}
For a transitive action $U\leq S_w$ write
$a(U)=|N_{S_w}(U)|$. If a profile contains $m_U$ copies of each action
type, then the number of labelled orbit data is
\begin{equation}\label{eq:weights}
 \frac{n!}{\prod_U a(U)^{m_U}m_U!}.
\end{equation}
Indeed, partition the points into the prescribed unordered blocks and
choose a conjugate of $U$ on each block. The factors $w!$ cancel.
A subgroup full on each block recovers both its actual orbits and their
permutation actions, so there is no further overcount.

For later use, let $b=n-w$. Pointing an original orbit of type $U$
has placement factor $n!/(b!a(U))$. If its complete complement is
$J\leq S_b$, a fibre bound $z_U(J)$ therefore gives the normalized
bound
\begin{equation}\label{eq:pointing}
 \frac{1}{b!a(U)p_nG_R}\sum_{J\leq S_b}z_U(J).
\end{equation}
In particular a uniform fibre envelope $z_U(b)$ produces the coefficient
\begin{equation}\label{eq:delta}
 \Delta(n,b)\frac{z_U(b)}{a(U)},\qquad
 \Delta(n,b)=\frac{p_bG_{\lfloor b/2\rfloor}}{p_nG_R}
\end{equation}
on the complete normalized count $s_b/L_b$.
An auxiliary quotient cover does not change $a(U)$.

\subsection{Algebraic inputs}
We use Gaussian coefficients, Goursat's lemma, Burnside's basis
theorem, the inflation--restriction sequence, the degree-two
K\"unneth formula, and classification of group extensions by $H^2$.
For permutation groups we use the abelian-section bound of
Kov\'acs and Praeger~\cite{KovacsPraeger1989}: an elementary abelian
$p$-section of a subgroup of $S_b$ has rank at most $b/p$.
Its form for $p$-groups is also recorded as
\cite[Theorem 2.8]{RoneyDougalTracey2025}.
The module-generator estimates used in Section~\ref{sec:capacity}
come from \cite[Definition 1.5, Theorem 1.6 and Corollary 3.12]{Tracey2018}.
The finite algebraic constants and the distinction between literal
witness verification and exhaustive coverage are specified in
Appendix~\ref{sec:certificates}.
